\documentclass[11pt]{article}
\usepackage[margin=1in]{geometry}
\usepackage{amsmath,amssymb,amsthm}
\usepackage{graphicx}
\usepackage[hidelinks]{hyperref}

\newtheorem{theorem}{Theorem}
\newtheorem{corollary}{Corollary}
\newtheorem{remark}{Remark}

\DeclareMathOperator{\Lip}{Lip}

\title{A Scaling Reduction for Equality in the Multiplicative Infinitesimal Prekopa--Leindler Inequality}
\author{fa\_banach\_001 lane 0}
\date{2026-07-19}

\begin{document}
\maketitle

\section*{Status}

\textbf{partial result / reduction}.  The result below does not characterize all
geometric equality cases.  It reduces the equality problem for the
multiplicative inequality to the equality problem for the additive mean-form
inequality, up to one scalar normalization.

\section*{Source Problem}

The source paper is Sotiris Armeniakos and Jacopo Ulivelli,
``A new infinitesimal form of the Prekopa--Leindler inequality with
multiplicative structure and applications'', arXiv:2602.10812.  After proving
Theorem 1.1 and Proposition 2.8, the authors state that a full characterization
of equality cases for the multiplicative form is outside the paper and left for
future work.

\begin{center}
\includegraphics[width=\textwidth]{figures/open_problem_crop.png}
\end{center}

The source defines three bilinear forms:
\[
\begin{aligned}
\langle \rho_0,\rho_1\rangle_P
&=
\int_{\partial K}\langle {\rm II}^{-1}\nabla_{\partial K}\rho_0,
\nabla_{\partial K}\rho_1\rangle\,d\mu
-\int_{\partial K}H_\mu\rho_0\rho_1\,d\mu \\
&\qquad
+\frac{1}{\mu(K)}
\left(\int_{\partial K}\rho_0\,d\mu\right)
\left(\int_{\partial K}\rho_1\,d\mu\right),\\
\langle \varphi_0,\varphi_1\rangle_{BL}
&=
\int_K\langle(\nabla^2u)^{-1}\nabla\varphi_0,\nabla\varphi_1\rangle\,d\mu
-\int_K\varphi_0\varphi_1\,d\mu\\
&\qquad
+\frac{1}{\mu(K)}
\left(\int_K\varphi_0\,d\mu\right)
\left(\int_K\varphi_1\,d\mu\right),\\
\langle \rho,\varphi\rangle_I
&=
\int_{\partial K}\rho\varphi\,d\mu
-\frac{1}{\mu(K)}
\left(\int_{\partial K}\rho\,d\mu\right)
\left(\int_K\varphi\,d\mu\right).
\end{aligned}
\]
The two inequalities are
\begin{align}
\langle \rho,\varphi\rangle_I
&\le
\frac{\langle\rho,\rho\rangle_P+\langle\varphi,\varphi\rangle_{BL}}{2},
\tag{M}\\
\langle\rho,\varphi\rangle_I^2
&\le
\langle\rho,\rho\rangle_P\langle\varphi,\varphi\rangle_{BL}.
\tag{P}
\end{align}

\section*{Proof Intuition}

The multiplicative inequality is obtained in the source by applying an
abstract quadratic discriminant argument to the mean inequality.  The same
quadratic keeps exact equality information.  For fixed \(\rho,\varphi\), apply
the mean inequality to \(t\rho\) and \(\varphi\).  This gives a nonnegative
quadratic in \(t\).  Multiplicative equality says precisely that this
quadratic has zero discriminant, hence that it touches zero at one scalar
normalization of \(\rho\).  That touching point is exactly equality in the
mean inequality.

\begin{theorem}[Scaling reduction for product equality]
Let \(K,u,\mu\) satisfy the hypotheses of Theorem 1.1 of arXiv:2602.10812, and
let \(\rho\) and \(\varphi\) be admissible test functions.  Put
\[
A=\langle\rho,\rho\rangle_P,\qquad
B=\langle\varphi,\varphi\rangle_{BL},\qquad
C=\langle\rho,\varphi\rangle_I .
\]
Assume \(A>0\) and \(B>0\).  Then equality holds in the multiplicative
inequality \((P)\) for \((\rho,\varphi)\) if and only if equality holds in the
mean inequality \((M)\) for \((\lambda\rho,\varphi)\), where
\[
  \lambda=\frac{C}{A}.
\]
Equivalently, every nontrivial product-equality pair is a scalar
re-normalization of a mean-equality pair, and conversely.
\end{theorem}

\begin{proof}
Apply \((M)\) to \(t\rho\) and \(\varphi\), for arbitrary \(t\in\mathbb R\).
By bilinearity and homogeneity, this is
\[
  tC\le \frac{t^2A+B}{2},
\]
or
\[
  q(t):=At^2-2Ct+B\ge 0
  \qquad\text{for every }t\in\mathbb R.
\]
Since \(A>0\), the quadratic \(q\) is nonnegative for all \(t\).

If equality holds in \((P)\), then \(C^2=AB\).  Hence
\[
  q(C/A)=A(C/A)^2-2C(C/A)+B=0.
\]
Thus equality holds in \((M)\) for \((\lambda\rho,\varphi)\) with
\(\lambda=C/A\).

Conversely, suppose equality holds in \((M)\) for \((\lambda\rho,\varphi)\)
for some real \(\lambda\).  Then \(q(\lambda)=0\).  Because \(q(t)\ge0\) for
all \(t\) and \(A>0\), a real zero forces the discriminant to vanish:
\[
  (-2C)^2-4AB=0.
\]
Therefore \(C^2=AB\), which is equality in \((P)\).
\end{proof}

\begin{remark}[Trivial equality cases]
The source records that \(\langle\rho,\rho\rangle_P=0\) if and only if
\(\rho\equiv0\), and the Brascamp--Lieb kernel on compact sets consists of
constants.  Thus \(\rho=0\) or constant \(\varphi\) gives the trivial equality
\(0=0\) in \((P)\).  The theorem above isolates all nontrivial cases.
\end{remark}

\begin{corollary}[Two-scale product equality family]
Let \((\rho_0,\varphi_0)\) be any nontrivial equality pair for the mean
inequality \((M)\); in particular, one may take the homothetic
Prekopa--Leindler family supplied by Proposition 2.8 of the source paper:
\[
  \rho_0=\alpha h_{K+x_0}(\nu_K),\qquad
  \varphi_0(x)=\alpha u^*(\nabla u(x-x_0))+z .
\]
Then for every real \(s,t\), and after adding an arbitrary constant to the
second component if desired, the pair
\[
  (s\rho_0,t\varphi_0)
\]
is an equality case for the multiplicative inequality \((P)\).  If
\((\rho_0,\varphi_0)\) is nontrivial, this rescaled pair is an equality case
for the mean inequality \((M)\) precisely when \(s=t\).
\end{corollary}

\begin{proof}
The first claim follows from the homogeneity
\[
\langle s\rho_0,t\varphi_0\rangle_I^2
=s^2t^2\langle\rho_0,\varphi_0\rangle_I^2
=s^2t^2
\langle\rho_0,\rho_0\rangle_P
\langle\varphi_0,\varphi_0\rangle_{BL}.
\]
This is exactly
\[
\langle s\rho_0,s\rho_0\rangle_P
\,
\langle t\varphi_0,t\varphi_0\rangle_{BL}
\]
by homogeneity of the two quadratic forms.  Constants may be added to
\(\varphi_0\) because both
\(\langle\cdot,\cdot\rangle_{BL}\) and \(\langle\cdot,\cdot\rangle_I\) vanish
on the constant part of the second variable.

For the final assertion, a nontrivial mean equality pair satisfies
\[
\langle\rho_0,\varphi_0\rangle_I
=
\langle\rho_0,\rho_0\rangle_P
=
\langle\varphi_0,\varphi_0\rangle_{BL}>0
\]
by the arithmetic--geometric mean equality condition.  Therefore mean equality
for \((s\rho_0,t\varphi_0)\) is equivalent to
\[
  st=\frac{s^2+t^2}{2},
\]
which is \((s-t)^2=0\).
\end{proof}

\section*{Limitations and Novelty Check}

This is a reduction, not the full equality-case classification.  It shows that
future work on the multiplicative equality cases can focus on the mean-form
equality cases: once those are known, all product equality cases follow by the
single scalar de-normalization above, together with the trivial
zero/constant-kernel cases.

The cheap run indexes had no prior entry for arXiv:2602.10812.  A web search
for the exact title and equality-case phrases found no later equality-case
resolution.  A same-author follow-up, arXiv:2602.23350, treats a strengthened
dimensional Brunn--Minkowski condition and the (B)-conjecture, not the equality
classification considered here.

\section*{References}

\begin{thebibliography}{9}
\bibitem{AU}
S. Armeniakos and J. Ulivelli,
\emph{A new infinitesimal form of the Prekopa--Leindler inequality with
multiplicative structure and applications},
arXiv:2602.10812.

\bibitem{AUfollow}
S. Armeniakos and J. Ulivelli,
\emph{A strengthening of the dimensional Brunn--Minkowski conjecture implies
the (B)-conjecture},
arXiv:2602.23350.
\end{thebibliography}

\end{document}
